% Noeris: arXiv submission draft (v1)
% Compile: pdflatex noeris-submission.tex && pdflatex noeris-submission.tex
%
% No external .sty required — uses standard article class with geometry.

\documentclass[11pt]{article}

\usepackage[margin=1in]{geometry}
\usepackage[utf8]{inputenc}
\usepackage[T1]{fontenc}
\usepackage{hyperref}
\usepackage{url}
\usepackage{booktabs}
\usepackage{amsfonts}
\usepackage{amsmath}
\usepackage{nicefrac}
\usepackage{microtype}
\usepackage{graphicx}
\usepackage{xcolor}
\usepackage{tikz}
\usetikzlibrary{shapes.geometric, arrows.meta, positioning, fit, backgrounds}
\usepackage{multirow}
\usepackage{enumitem}
\usepackage{natbib}

\hypersetup{
  colorlinks=true,
  linkcolor=blue!60!black,
  citecolor=blue!60!black,
  urlcolor=blue!60!black,
}

\title{Noeris: Architecture-Agnostic Kernel Fusion and Autotuning}

\author{
  Doruk Tan Ozturk \\
  Independent Researcher \\
  \texttt{github.com/PwnKit-Labs/noeris}
}

\date{}

\begin{document}

\maketitle

% ---------------------------------------------------------------------------
% ABSTRACT
% ---------------------------------------------------------------------------
\begin{abstract}
We present Noeris, an autonomous GPU kernel optimization system with three
distinguishing properties: \textbf{(1)}~parameterized kernel templates (no
free-form source rewriting), \textbf{(2)}~a shape-indexed cross-run
configuration database keyed by \texttt{(operator, shape\_bucket, hardware)}
that persists winning configurations across sessions, and
\textbf{(3)}~complementary selectors---a gradient-boosted cost model and a
multi-armed bandit---that rank candidates before GPU evaluation.

A fused Triton kernel for the Gemma~3/4 attention prologue
(\texttt{Q-RMSNorm $\to$ K-RMSNorm $\to$ Q-RoPE $\to$ K-RoPE}) achieves
$10.2$--$12.9\times$ speedup on A100 and $10.4$--$11.9\times$ on H100,
reaching 1627.7~GB/s (${\approx}49\%$ of HBM3 peak). Cross-model
generalization across 19 models from 13 families confirms architecture-general
fusion ($6.5$--$8.9\times$ on T4, $3.9$--$9.8\times$ on A100). A
sliding-window attention kernel with tile-pruning beats cuDNN FlashAttention
by up to $3.56\times$ on T4 and $6.24\times$ on A100 (8/8 shapes win). A
26-layer Gemma~4 forward pass achieves $1.18\times$ end-to-end speedup on T4
and $1.43\times$ on A100 by fusing only operations accounting for ${<}20\%$ of
per-layer compute. Cross-hardware fused kernel rankings achieve Spearman
$\rho = 0.907$ ($p < 10^{-7}$), validating zero-shot config prediction. A
bandit convergence study reaches ${\geq}98\%$ of exhaustive-search optimal in
one iteration (12\% of the search space). The system covers 21 operators
(18/20 validated on T4), 110+ shape buckets, and 606 unit tests at
${\sim}\$0.01$ per iteration. Compiler analysis shows \texttt{torch.compile}
reduces launches $40 \to 9$; Noeris: $40 \to 1$. Code and data:
\url{https://github.com/PwnKit-Labs/noeris} (MIT).
\end{abstract}

% ---------------------------------------------------------------------------
% 1. INTRODUCTION
% ---------------------------------------------------------------------------
\section{Introduction}
\label{sec:intro}

LLM-driven GPU kernel optimization has become an active research topic in
2025--2026~\citep{autokernel,kernelskill,cudal1,cudaagent,kernelfoundry}. These
systems share a common architecture: an LLM agent proposes kernel code, a
harness measures correctness and performance, and an orchestration layer decides
whether to keep or revert the change.

A shared limitation of published systems is that \textbf{search state does not
persist across sessions}. Each invocation starts from the same initial kernel (or
a cached version), runs its own iterative search loop, and discards the
trajectory. Information learned in one run---which tile sizes work for a given
matrix shape on a given GPU---is not systematically reused in the next.

Noeris investigates an alternative. Rather than rewriting kernel source per
invocation, we generate Triton~\citep{triton} kernels from a compact parameter tuple (e.g.,
\texttt{BLOCK\_SIZE\_M}, \texttt{BLOCK\_SIZE\_N}, \texttt{num\_warps},
\texttt{num\_stages}) and store winning configurations in a \textbf{shape-indexed
cross-run database} keyed by \texttt{(operator, shape\_bucket, hardware)}. When
an LLM proposer is invoked, it sees the database state as cross-run insights,
allowing it to reason about what has worked on similar shapes before. A
gradient-boosted cost model trained on accumulated benchmark data then
rank-orders grid candidates at prediction time, without incurring additional GPU
calls.

All code, benchmark data, and reproduction scripts are available at
\url{https://github.com/PwnKit-Labs/noeris} under the MIT License. We make thirteen contributions:

\begin{enumerate}[leftmargin=*,itemsep=2pt]

\item \textbf{A practical fused kernel that makes QK-norm+RoPE fusion actually
work.} The Gemma~3/4 attention prologue
(\texttt{Q-RMSNorm}$\to$\texttt{K-RMSNorm}$\to$\texttt{Q-RoPE}$\to$\texttt{K-RoPE})
has prior art: vLLM has an experimental \texttt{enable\_qk\_norm\_rope\_fusion}
pass (disabled by default due to H100 performance regression), SGLang has
\texttt{fused\_qknorm\_rope}, and Liao et al.\ describe algebraic optimizations
in ``Flash Normalization.'' Our parameterized Triton kernel with bandit-tuned
configs beats the separated-kernel-launches baseline by 10.2--12.9$\times$ on
A100 and 10.4--11.9$\times$ on H100 across six Gemma~3/4 shape buckets (60/60
correct, zero failures). Peak 1627.7~GB/s on H100.

\item \textbf{System.} A complete autonomous kernel search loop for 21
parameterized Triton operators (18/20 validated on T4), running on cloud GPUs
via Modal with GitHub Actions orchestration, at ${\sim}\$0.01$ per iteration.

\item \textbf{Full Gemma~4 attention support.} Grouped-query attention, asymmetric
local/global \texttt{head\_dim} (256/512), and Gemma-mode $(1+w)$ RMSNorm affine.

\item \textbf{Learned feasibility.} Configurations that fail at runtime flow
into the bandit as reward=0 samples, eliminating all hand-coded shared-memory
filters.

\item \textbf{Evaluation.} First direct reproduction of AutoKernel's H100
memory-bound kernel numbers, with substantial improvements: 11.66$\times$
RMSNorm (+120\%), 9.65$\times$ Cross-entropy (+228\%), 6.38$\times$ Softmax
(+85\%).

\item \textbf{Learned cost model + complementary selectors + adaptive routing.}
GBR cost model (R$^2$=0.94); three-way selector comparison shows cost model and
bandit are complementary; adaptive router matches best fixed selector within
0.5\%.

\item \textbf{Compiler failure analysis (\S\ref{sec:compiler-analysis}).}
\texttt{torch.compile} reduces launches 40$\to$9 but still generates 4 separate
Triton kernels. Noeris: 1 kernel, 1 launch, 6.08$\times$ vs.\ eager, 1.62$\times$ vs.\
compiled.

\item \textbf{Cross-model generalization (\S\ref{sec:cross-model}).} 19 model
shapes, 13 families, all achieve 6.5--8.9$\times$ fusion speedup on T4.

\item \textbf{End-to-end 26-layer benchmark (\S\ref{sec:e2e}).} 1.18$\times$ on
T4, 1.43$\times$ on A100---6--21$\times$ more E2E impact than vLLM's 2--3\%.

\item \textbf{Sliding-window attention beats cuDNN
(\S\ref{sec:swa-eval}).} Up to 3.56$\times$ on T4, 6.24$\times$ on A100 via
tile-pruning.

\item \textbf{Cross-hardware transfer (\S\ref{sec:cross-hw-transfer}).}
Spearman $\rho=0.907$ ($p<10^{-7}$) between A100 and H100 fused kernel latency
rankings, validating zero-shot config prediction.

\item \textbf{Bandit convergence (\S\ref{sec:bandit-convergence}).} $\geq$98\%
of exhaustive-search optimal in one iteration with 12\% of the search space.

\item \textbf{Mamba-3 SSM scan operator.} Validates architecture-agnostic
coverage beyond transformer attention.

\end{enumerate}


% ---------------------------------------------------------------------------
% 2. SYSTEM
% ---------------------------------------------------------------------------
\section{System}
\label{sec:system}

% TikZ system architecture diagram
\begin{figure}[t]
\centering
\begin{tikzpicture}[
    node distance=1.2cm and 1.0cm,
    box/.style={
      rectangle, rounded corners=3pt, draw, thick,
      minimum width=2.8cm, minimum height=1.2cm,
      text width=2.6cm, align=center, font=\small
    },
    bluebox/.style={box, fill=blue!12, draw=blue!50!black},
    greenbox/.style={box, fill=green!12, draw=green!50!black,
      minimum width=10.2cm, text width=10.0cm},
    graybox/.style={box, fill=gray!15, draw=gray!60!black},
    arr/.style={-{Stealth[length=3mm]}, thick},
    lbl/.style={font=\scriptsize\itshape, fill=white, inner sep=1pt},
  ]

  % Top row: three processing stages
  \node[bluebox] (proposer)
    {LLM\\Proposer\\{\scriptsize + cross-run insights}};
  \node[bluebox, right=of proposer] (costmodel)
    {Cost Model\\(GBR)};
  \node[bluebox, right=of costmodel] (bandit)
    {Bandit /\\Adaptive Router};

  % Middle: database
  \node[greenbox, below=1.4cm of costmodel] (db)
    {Shape-Indexed Config Database\\
     {\scriptsize key: (operator, shape\_bucket, hardware) \quad
      value: best\_config, metric, trajectory}};

  % Bottom: GPU runner
  \node[graybox, below=1.2cm of db] (runner)
    {Modal GPU Runner\\(A100 / H100)\\{\scriptsize \raisebox{0pt}{\~{}}\$0.01\,/\,iteration}};

  % Arrows: top row
  \draw[arr] (proposer) -- node[lbl, above] {candidates} (costmodel);
  \draw[arr] (costmodel) -- node[lbl, above] {rank-orders} (bandit);

  % Arrows: top row to/from database
  \draw[arr] (bandit.south) -- node[lbl, right, pos=0.4] {configs} (db.north -| bandit);
  \draw[arr] (db.north -| proposer) -- node[lbl, left, pos=0.4] {history} (proposer.south);

  % Arrows: database <-> runner
  \draw[arr] (db.south) -- node[lbl, right, pos=0.45] {configs} (runner.north);
  \draw[arr] ([xshift=-1.2cm]runner.north) -- node[lbl, left, pos=0.45] {results} ([xshift=-1.2cm]db.south);

  % Dashed outer box
  \begin{scope}[on background layer]
    \node[draw=gray!70, dashed, thick, rounded corners=5pt,
          fit=(proposer)(bandit)(db)(runner),
          inner xsep=12pt, inner ysep=10pt,
          label={[font=\small\bfseries, anchor=north]above:Noeris}] {};
  \end{scope}
\end{tikzpicture}
\caption{Noeris system architecture. The LLM proposer generates configurations
informed by cross-run database history. A gradient-boosted cost model and a
multi-armed bandit (or adaptive router combining both) filter candidates before
GPU evaluation. Results flow back into the database, compounding across CI runs.}
\label{fig:system}
\end{figure}


\subsection{Parameterized Kernels and the Operator Registry}
\label{sec:operator-registry}

Every operator in Noeris is described by a \texttt{TritonOperatorSpec}:

\begin{quote}\small
\texttt{name} (e.g.\ matmul, rmsnorm),
\texttt{param\_space} (e.g.\ \texttt{\{BLOCK\_SIZE: [128, 256, \ldots]\}}),
\texttt{curated\_configs} (hand-picked starting points),
\texttt{shape\_buckets} (workload-shaped test shapes),
\texttt{metric\_name} (\texttt{tflops} or \texttt{gb\_per\_s}),
\texttt{grid\_generator\_fn} (systematic parameter grid).
\end{quote}

Each operator is registered once and dispatched uniformly by the rest of the
system: the proposer, selector, database, and runner are all operator-agnostic.
Adding a new operator requires approximately 200 lines plus a registry entry.

The parameterized interface is a deliberate constraint. By prohibiting free-form
source rewriting, we guarantee that every benchmark result can be indexed by a
stable, typed key and that cost model feature extraction is deterministic. The
tradeoff---reduced expressiveness relative to free-form source agents---is
accepted in exchange for tractability.


\subsection{Shape-Indexed Cross-Run Config Database}
\label{sec:config-db}

\texttt{ConfigDatabase} persists benchmark results as JSON, keyed by
\texttt{\{operator\}:\{shape\_bucket\}:\{hardware\}}. Shape buckets are
operator-specific: for matmul we classify by $(M \times N \times K$, aspect
ratio, $K$ ratio); for rmsnorm by \texttt{hidden\_dim}; for attention by
(\texttt{seq\_len}, \texttt{head\_dim}). Configurations learned on
\texttt{llama\_7b} shapes are not reused for \texttt{mixtral} shapes---each
bucket maintains its own incumbent.

Across CI runs, the database is restored from the previous successful artifact,
updated with new results, and saved as a new artifact. This is a standing JSON
database that compounds knowledge over time without requiring a persistent
server process.


\subsection{LLM Proposer with Cross-Run Insights}
\label{sec:llm-proposer}

When the LLM proposer runs, it receives: (1)~the operator's parameter space and
hardware constraints, (2)~the target shapes for the current iteration, and
(3)~cross-run insights extracted from the database: per bucket, the top-3
configurations and their best measured metrics. The proposer responds with up to
4 novel configurations plus a natural-language rationale.


\subsection{Frontier-Aware Config Selection}
\label{sec:frontier}

\texttt{select\_configs\_for\_operator} allocates up to $N$ slots per iteration
with explicit semantics: (1)~\textbf{Incumbent}---the best known config;
(2)~\textbf{LLM-proposed}---novel configs from the proposer;
(3)~\textbf{Curated}---hand-picked starters not yet tested;
(4)~\textbf{Exploration}---systematic grid configs not yet tested. This ensures
known-good configurations are always re-validated while allocating explicit
budget to exploration and novelty.


\subsection{Learned Cost Model}
\label{sec:cost-model}

The central technical contribution beyond the LLM proposer is a
\textbf{learned cost model} trained on the shape-indexed database. Features are a
fixed-width 20-dimensional vector including: operator ID (one-hot), hardware ID
(one-hot), shape dimensions (5 slots), configuration parameters (8 slots), and
derived features (\texttt{log(shape\_product)}, \texttt{log(max\_dim)},
\texttt{log(tile\_area)}, etc.).

We use sklearn \texttt{GradientBoostingRegressor} with 200 trees, depth~5,
learning rate 0.05. On an 80/20 split of 384 training points, the model achieves
$R^2 = 0.970$. Feature importance analysis reveals that
\texttt{log\_shape\_min} (0.433) and \texttt{log\_shape\_max} (0.126) are
dominant, followed by raw shape dimensions and \texttt{num\_warps}.

At selection time, the selector gathers up to 40 grid candidates and calls
\texttt{cost\_model.rank\_configs()}, returning them sorted by predicted metric.
The cost model is deterministic: given the same training data, it always
produces the same ranking.


\subsection{Execution Backend (Modal)}
\label{sec:modal}

Benchmark scripts are self-contained Python files. A single Modal function
executes them on a warm GPU container. For multi-iteration workflows, a
persistent session context keeps one container warm across all iterations,
cutting per-call overhead from ${\sim}10$\,s to 1--3\,s. Total cost to reproduce
all results: approximately \$1.44 in Modal GPU credits.


% ---------------------------------------------------------------------------
% 3. OPERATORS
% ---------------------------------------------------------------------------
\section{Operators}
\label{sec:operators}

\begin{table}[t]
\centering
\caption{Noeris operator registry (21 operators, 110+ shape buckets).}
\label{tab:operators}
\small
\begin{tabular}{@{}llcr@{}}
\toprule
Operator & Parameter space & Metric & Buckets \\
\midrule
matmul & \texttt{BLOCK\_M/N/K, GROUP\_SIZE\_M, warps, stages} & TFLOPS & 10 \\
rmsnorm & \texttt{BLOCK\_SIZE, warps, stages} & GB/s & 12 \\
softmax & \texttt{BLOCK\_SIZE, warps, stages} & GB/s & 7 \\
layernorm & \texttt{BLOCK\_SIZE, warps, stages} & GB/s & 8 \\
cross\_entropy & \texttt{BLOCK\_SIZE, warps, stages} & GB/s & 9 \\
attention (FA+SWA) & \texttt{BLOCK\_M, BLOCK\_N, warps, stages, \ldots} & TFLOPS & 12 \\
rotary (RoPE) & \texttt{BLOCK\_SIZE, warps, stages} & GB/s & 8 \\
geglu & \texttt{BLOCK\_SIZE, warps, stages} & GB/s & 5 \\
qk\_norm\_rope (fwd) & \texttt{BLOCK\_SIZE, warps, stages} & GB/s & 6 \\
qk\_norm\_rope (bwd) & \texttt{BLOCK\_SIZE, warps, stages} & GB/s & 6 \\
moe\_router & \texttt{BLOCK\_SIZE, warps, stages} & GB/s & 4 \\
grouped\_gemm & \texttt{BLOCK\_M/N/K, warps, stages} & TFLOPS & 4 \\
paged\_kv\_decode & \texttt{BLOCK\_SIZE, warps, stages} & GB/s & 4 \\
ple\_gather & \texttt{BLOCK\_SIZE, warps, stages} & GB/s & 2 \\
ssm\_scan (Mamba-3) & \texttt{BLOCK\_B/D/L, warps, stages} & GB/s & 3 \\
silu\_mul & \texttt{BLOCK\_SIZE, warps, stages} & GB/s & 4 \\
embedding & \texttt{BLOCK\_SIZE, warps, stages} & GB/s & 3 \\
add\_residual & \texttt{BLOCK\_SIZE, warps, stages} & GB/s & 3 \\
dropout & \texttt{BLOCK\_SIZE, warps, stages} & GB/s & 3 \\
fused\_bias\_act & \texttt{BLOCK\_SIZE, warps, stages} & GB/s & 4 \\
quantize\_fp8 & \texttt{BLOCK\_SIZE, warps, stages} & GB/s & 3 \\
\bottomrule
\end{tabular}
\end{table}


\subsection{Sliding-Window Attention Kernel}
\label{sec:swa-kernel}

The attention operator supports a \texttt{WINDOW\_SIZE} constexpr that restricts
each query token to at most \texttt{WINDOW\_SIZE} key tokens, enabling tile
pruning for local attention patterns. When \texttt{WINDOW\_SIZE > 0}, the inner
K-tile loop is bounded by:
\[
k_\text{start} = \max(0,\; \mathit{pid} \cdot \texttt{BLOCK\_M} - \texttt{WINDOW\_SIZE} + 1), \quad
k_\text{end} = \min(N,\; \mathit{pid} \cdot \texttt{BLOCK\_M} + \texttt{WINDOW\_SIZE})
\]
rather than iterating over all K tiles from 0 to $N$. For a representative
Gemma~4 configuration---\texttt{window=1024}, \texttt{seq=4096},
\texttt{BLOCK\_N=64}---this prunes approximately half the K tiles.


\subsection{Fused QK-RMSNorm + RoPE Prologue (Headline Result)}
\label{sec:fused-prologue}

The Gemma~3/4 attention prologue
(\texttt{Q-RMSNorm}$\to$\texttt{K-RMSNorm}$\to$\texttt{Q-RoPE}$\to$\texttt{K-RoPE})
is a known fusion target. vLLM~\citep{vllm} has an experimental
\texttt{enable\_qk\_norm\_rope\_fusion} pass---a \texttt{torch.compile} pass with a
CUDA kernel---but it is \textbf{disabled by default} due to performance regression
on H100. The reported benefit when enabled is 2--3\% end-to-end. Our
contribution is a parameterized Triton implementation with bandit-tuned
configurations that makes this fusion practical and performant across hardware.

\texttt{qk\_norm\_rope\_kernel} is a fused Triton kernel that for each row in a
$(B, H, S, D)$ tensor: (i)~loads the row, (ii)~computes
$\mathit{rstd} = \text{rsqrt}(\text{mean}(x^2) + \epsilon)$, (iii)~applies the
Gemma-mode $(1+w)$ affine, (iv)~rotates the resulting even/odd pairs with the
RoPE $\cos$/$\sin$ tables, and (v)~stores the result. Two launches (one Q, one
K) replace the four of the separated path.

\textbf{Measurement.} 6 Gemma~3/4 shape buckets $\times$ 5 curated configs
$\times$ 2 GPUs (A100-SXM4-40GB and H100-SXM5-80GB). Timer:
\texttt{cuda\_event} + L2 flush, 3 warmup / 10 trials, median. Correctness:
$\texttt{max\_err} \leq 0.1$ against PyTorch fp32 reference. \textbf{All 60
(shape, config) combinations are correct---zero failures.}

\begin{table}[t]
\centering
\caption{Fused QK-RMSNorm+RoPE: best fusion speedup per Gemma shape.}
\label{tab:fusion-speedup}
\small
\begin{tabular}{@{}lcccccc@{}}
\toprule
Shape & GQA & head\_dim & A100 GB/s & A100 speedup & H100 GB/s & H100 speedup \\
\midrule
gemma4\_31b\_global & 32:4 & 512 & 925.5 & \textbf{12.85$\times$} & \textbf{1627.7} & 11.88$\times$ \\
gemma4\_26b\_a4b\_global & 16:2 & 512 & 905.0 & 12.40$\times$ & 1576.5 & 11.73$\times$ \\
gemma4\_31b\_local & 32:16 & 256 & 731.2 & 11.30$\times$ & 1536.3 & 11.81$\times$ \\
gemma3\_local\_1024 & 16:16 & 256 & 715.2 & 10.69$\times$ & 1490.2 & \textbf{11.82$\times$} \\
gemma4\_26b\_a4b\_local & 16:8 & 256 & 711.1 & 10.37$\times$ & 1443.2 & 11.49$\times$ \\
gemma4\_e2b\_local & 8:1 & 256 & 593.6 & 10.23$\times$ & 1213.7 & 10.35$\times$ \\
\bottomrule
\end{tabular}
\end{table}

\textbf{Why so much higher than a pure HBM-accounting estimate?} A purely
HBM-accounting cost model predicts ${\sim}2\times$ fusion speedup. The measured
10--13$\times$ is 5--6$\times$ higher. The gap is \textbf{kernel launch overhead}.
At these tile sizes each individual kernel takes 30--100~$\mu$s on A100/H100;
CUDA launch latency (${\sim}$5--10~$\mu$s per launch) is 10--20\% of each
separated call, and going from 4 launches to 2 saves a disproportionate fraction
of wall-clock time. This is precisely the regime where Triton fusion provides
asymmetric value---not at the HBM-bandwidth frontier, but at the
launch-amortization frontier.

\textbf{Backward pass.} The backward kernel fuses the reverse sequence: RoPE
inverse rotation, Gemma-mode $(1+w)$ affine backward, and RMSNorm backward,
replacing 4 separate backward launches with 2 fused kernels. Validated on T4:
4.9--7.5$\times$ backward fusion speedup. To our knowledge, \textbf{no existing
framework fuses the combined QK-RMSNorm+RoPE backward into a single kernel}---
Liger Kernel fuses them separately.


\subsection{Mamba-3 SSM Scan Operator}
\label{sec:ssm-scan}

To demonstrate that the parameterized template approach extends beyond
transformer attention, we implement a selective state-space model (SSM) scan
operator targeting Mamba-3-style architectures~\citep{mamba2}. The operator
implements the parallel-scan recurrence ($h_t = A_t \cdot h_{t-1} + B_t \cdot
x_t$; $y_t = C_t \cdot h_t$) via the Blelloch prefix-sum algorithm, tiled across
batch and channel dimensions. The operator is registered through the standard
\texttt{TritonOperatorSpec} interface; no changes to the bandit, cost model, or
config database were required. GPU evaluation is pending.


% ---------------------------------------------------------------------------
% 4. EVALUATION
% ---------------------------------------------------------------------------
\section{Evaluation}
\label{sec:eval}

Every number is drawn directly from raw benchmark artifacts. No smoothing or
post-hoc selection has been applied.

\textbf{Transparency note.} Sections~\ref{sec:aggregate}--\ref{sec:autokernel}
use our \textbf{internal LLM-shape benchmark}: 53 problems with shapes drawn
from real LLM workloads, evaluated in FP16, using \texttt{triton.testing.do\_bench}.
Starting from \S\ref{sec:compiler-analysis}, results use the upstream KernelBench
methodology (\texttt{cuda\_event} + L2 flush, FP32 \texttt{nn.Module} problems).
The two regimes measure \textbf{different workloads}; both are reported.


\subsection{Experimental Setup}
\label{sec:setup}

\textbf{Hardware.} A100-SXM4-40GB (Modal), H100-SXM5-80GB (Modal), Tesla T4
(Kaggle/Colab free tier). 18 of 20 non-matmul operators validated on T4; all 21
operators validated on A100 and H100.

\textbf{Problem set.} 53 curated problems at Levels~1--3 from
KernelBench~\citep{kernelbench}: matmul (19), rmsnorm (7), softmax (8),
layernorm (6), cross\_entropy (7), attention (6). All FP16.

\textbf{Configs per problem.} 4 curated starters; best result taken. No search
iterations.

\textbf{Cost.} Full two-GPU evaluation: ${\sim}\$1.44$ in Modal credits.


\subsection{Aggregate Results}
\label{sec:aggregate}

\begin{table}[t]
\centering
\caption{fast$_p$ scores (fraction of problems with speedup $\geq t\times$ over PyTorch eager).}
\label{tab:fastp}
\small
\begin{tabular}{@{}lcccc@{}}
\toprule
 & \multicolumn{2}{c}{A100 (53 problems)} & \multicolumn{2}{c}{H100 (53 problems)} \\
\cmidrule(lr){2-3}\cmidrule(lr){4-5}
Threshold & Overall & Level 1/2 & Overall & Level 1/2 \\
\midrule
fast$_{1.0}$ & \textbf{56.6\%} & 58.3\% / 61.5\% & \textbf{56.6\%} & 58.3\% / 61.5\% \\
fast$_{1.5}$ & 41.5\% & 45.8\% / 42.3\% & 43.4\% & 45.8\% / 46.2\% \\
fast$_{2.0}$ & 37.7\% & 37.5\% / 42.3\% & 41.5\% & 45.8\% / 42.3\% \\
fast$_{3.0}$ & 32.1\% & 33.3\% / 34.6\% & 30.2\% & 33.3\% / 30.8\% \\
\bottomrule
\end{tabular}
\end{table}


\subsection{Per-Operator Deep Dive}
\label{sec:per-operator}

\textbf{RMSNorm} is the clearest win. On H100, the Mixtral shape reaches
\textbf{2625.1~GB/s at 11.66$\times$} PyTorch eager; LLaMA shapes achieve
10--11$\times$. The Triton kernel fuses squared-mean reduction, reciprocal-sqrt,
and element-wise scale into a single pass.

\textbf{Softmax} shows the largest shape-sensitivity. On H100, speedups range
from 2.18$\times$ (tiny) to \textbf{6.38$\times$} (vocab\_llama, width=32000).

\textbf{Cross-Entropy} delivers best absolute numbers: 2407.1~GB/s at
\textbf{9.65$\times$} on H100.

\textbf{LayerNorm} is weakest: 1.25--1.53$\times$ on H100. PyTorch's fused ATen
kernel is competitive.

\textbf{Matmul}: cuBLAS is hard to beat. 0.66--1.01$\times$ range. The two
LLaMA-7B shapes reach 1.01$\times$ on H100, tying cuBLAS at 692~TFLOPS.

\textbf{GeGLU}: all four Gemma shapes achieve $\geq$3$\times$ on A100 (3.08--3.98$\times$),
reaching 60--68\% of H100 peak HBM3 bandwidth with no search iterations.


\subsection{Comparison to AutoKernel}
\label{sec:autokernel}

\begin{table}[t]
\centering
\caption{Best-shape comparison to AutoKernel~\citep{autokernel} on H100.}
\label{tab:autokernel}
\small
\begin{tabular}{@{}lccc@{}}
\toprule
Kernel & Noeris best & AutoKernel & Delta \\
\midrule
RMSNorm & \textbf{11.66$\times$} (2625 GB/s) & 5.29$\times$ & \textbf{+120\%} \\
Cross-entropy & \textbf{9.65$\times$} (2407 GB/s) & 2.94$\times$ & \textbf{+228\%} \\
Softmax & \textbf{6.38$\times$} (2526 GB/s) & 3.44$\times$ & \textbf{+85\%} \\
LayerNorm & 1.53$\times$ (1667 GB/s) & 3.21$\times$\textsuperscript{$\dagger$} & $-$52\% \\
\bottomrule
\multicolumn{4}{@{}l}{\footnotesize $\dagger$ AutoKernel's LayerNorm baseline is \texttt{torch.compile}, not eager.}
\end{tabular}
\end{table}


\subsection{Cross-Run Learning Ablation (Honest Negative Result)}
\label{sec:cross-run-ablation}

At 5 iterations with strong curated priors, database-guided LLM proposals do not
outperform stateless proposals within the measurement noise floor. Matmul:
$-2.16\%$; rmsnorm: $-1.25\%$---neither statistically significant, both within
the ${\sim}$2--3\% GPU-runner noise floor. The negative result reflects the
strength of curated starters, not evidence against persistent learning in
principle.


\subsection{Cost Model Validation}
\label{sec:cost-model-validation}

With curated configs disabled (\texttt{--no-curated}), forcing both conditions to
start from the raw parameter grid:

\begin{table}[t]
\centering
\caption{Cost model ablation: filtered vs.\ unfiltered grid search (A100).}
\label{tab:cost-model}
\small
\begin{tabular}{@{}lccc@{}}
\toprule
Operator & Baseline (GB/s) & Filtered (GB/s) & $\Delta$ \\
\midrule
cross\_entropy & 682.16 & \textbf{936.94} & \textbf{+37.35\%} \\
softmax & 948.35 & \textbf{998.26} & +5.26\% \\
rmsnorm & 801.83 & 803.16 & +0.17\% \\
layernorm & 781.45 & 782.62 & +0.15\% \\
\bottomrule
\end{tabular}
\end{table}

The gain scales with grid size / budget: operators where exhaustive evaluation
would be expensive benefit most. For cross\_entropy, the filtered condition
reaches its near-final value (935.8~GB/s) on \textbf{iteration~1}---before the
baseline has found anything above 682~GB/s.


\subsection{Complementary Selectors: Bandit Fills the Sparse-Data Gap}
\label{sec:selectors}

\begin{table}[t]
\centering
\caption{Three-way selector comparison (\texttt{--no-curated}, 5 iterations $\times$ 6 configs, A100).}
\label{tab:selectors}
\small
\begin{tabular}{@{}lcccccc@{}}
\toprule
Operator & Baseline & Cost Model & $\Delta$ & Bandit & $\Delta$ & Winner \\
\midrule
matmul & 59.30 & 86.15 & +45\% & \textbf{138.38} & \textbf{+133\%} & Bandit \\
attention & 85.03 & \textbf{141.30} & \textbf{+66\%} & 134.95 & +59\% & Cost Model \\
cross\_entropy & 681.77 & \textbf{937.90} & +38\% & 937.05 & +37\% & Tied \\
softmax & 967.95 & \textbf{1018.50} & +5\% & 980.61 & +1\% & Cost Model \\
rmsnorm & 889.76 & 889.67 & 0\% & 883.57 & $-$1\% & No effect \\
\bottomrule
\end{tabular}
\end{table}

The two selectors occupy different niches. The bandit dominates matmul (+133\%
vs.\ +45\%) where training data is sparse and the cost model must extrapolate.
The cost model dominates attention (+66\% vs.\ +59\%) where the
\texttt{BLOCK\_M}$\times$\texttt{BLOCK\_N} product has a smooth, near-monotone
relationship with TFLOPS.


\subsection{Ensemble Failure}
\label{sec:ensemble-failure}

Naive alternation (half slots from cost model, half from bandit) matches the cost
model rather than the bandit on matmul (85.13 vs.\ 137.51~TFLOPS). The ensemble
wastes half the budget on the weaker selector's proposals when selectors
diverge.


\subsection{Hardware Cross-Learning}
\label{sec:hw-cross-learning}

\begin{table}[t]
\centering
\caption{A100$\to$H100 cost model transfer (memory-bound operators).}
\label{tab:hw-transfer}
\small
\begin{tabular}{@{}lcccc@{}}
\toprule
Operator & R$^2$ (cross-hw) & Spearman $\rho$ & Top-5 agree & Random \\
\midrule
rmsnorm & $-$0.108 & 0.961 & 60\% & 7.8\% \\
softmax & +0.463 & 0.988 & 60\% & 7.8\% \\
layernorm & $-$0.040 & 0.942 & 40\% & 7.8\% \\
cross\_entropy & +0.503 & 0.978 & 40\% & 7.8\% \\
\midrule
\textbf{Mean} & \textbf{+0.205} & \textbf{0.967} & \textbf{50\%} & \textbf{7.8\%} \\
\bottomrule
\end{tabular}
\end{table}

Rankings transfer perfectly (Spearman $\rho = 0.967$); absolute predictions do
not (R$^2 = 0.205$). The relative ordering of configurations is governed by the
same structural constraints on both GPUs---shared memory capacity, roofline
shape---while absolute throughput scales with bandwidth. Top-5 agreement is
50\% vs.\ 7.8\% random: a 6.4$\times$ lift.


\subsection{Multi-Trial Adaptive Router Validation}
\label{sec:adaptive-router}

\begin{table}[t]
\centering
\caption{Adaptive router: matmul A100, 3 trials $\times$ 5 iterations.}
\label{tab:router}
\small
\begin{tabular}{@{}lcccc@{}}
\toprule
Condition & Trial~1 & Trial~2 & Trial~3 & Mean $\pm$ $\sigma$ \\
\midrule
baseline & 58.33 & 59.01 & 54.92 & 57.42 $\pm$ 2.19 \\
cost\_model & 84.78 & 85.34 & 78.44 & 82.85 $\pm$ 3.83 \\
bandit & 136.07 & 137.51 & 124.85 & \textbf{132.81 $\pm$ 6.93} \\
naive\_ensemble & 85.33 & 85.60 & 78.22 & 83.05 $\pm$ 4.19 \\
adaptive\_router & 135.35 & 136.94 & 124.29 & \textbf{132.19 $\pm$ 6.89} \\
\bottomrule
\end{tabular}
\end{table}

The adaptive router matches the bandit within 0.5\% (132.19 vs.\ 132.81
TFLOPS) while the naive ensemble remains stuck at cost-model level (83.05). The
router converges to preferring the bandit arm within 1--2 iterations.


% ---------------------------------------------------------------------------
% 4.16 -- KEY RESULTS
% ---------------------------------------------------------------------------

\subsection{Compiler Failure Analysis: Why torch.compile Falls Short}
\label{sec:compiler-analysis}

\begin{table}[t]
\centering
\caption{Compiler comparison on T4: QK-RMSNorm+RoPE prologue.}
\label{tab:compiler}
\small
\begin{tabular}{@{}lccc@{}}
\toprule
Method & Kernels & Launches & Prologue (ms) \\
\midrule
PyTorch eager & 10 & 40 & 3.45 \\
\texttt{torch.compile} & 4 Triton + 5 DtoD & 9 & 0.92 \\
Noeris (fused) & 1 Triton & 1 & 0.57 \\
\midrule
\multicolumn{3}{@{}l}{Noeris speedup vs.\ eager} & \textbf{6.08$\times$} \\
\multicolumn{3}{@{}l}{Noeris speedup vs.\ compiled} & \textbf{1.62$\times$} \\
\bottomrule
\end{tabular}
\end{table}

The \textbf{40 $\to$ 9 $\to$ 1} launch progression tells the story.
\texttt{torch.\_dynamo.explain} reports 0 graph breaks, 1 graph, 30 ops---the
compiler has full visibility. Yet it \emph{chooses} not to fuse across the
RMSNorm reduction boundary. The generated kernel names
(\texttt{triton\_per\_fused\_mean\_pow} and
\texttt{triton\_poi\_fused\_rsqrt\_slice\_stack}) confirm HBM materialization
between each reduction and pointwise kernel pair.

At the full-layer level, Noeris achieves 1.54$\times$ over separated PyTorch
vs.\ the compiler's 1.23$\times$---a 25\% gap. Wrapping Noeris custom ops in
\texttt{torch.compile} is \emph{counterproductive} (1.20$\times$, slower than
Noeris alone) due to tracing overhead around opaque Triton ops.


\subsection{Cross-Model Generalization (19 Models, 13 Families)}
\label{sec:cross-model}

\begin{table}[t]
\centering
\caption{Cross-model prologue fusion on T4: 19 models, 13 families. All 19/19 pass correctness.}
\label{tab:cross-model-t4}
\small
\begin{tabular}{@{}lccc@{}}
\toprule
Model & QK-Norm? & Speedup & GB/s \\
\midrule
Qwen 3 32B & Yes & \textbf{8.90$\times$} & \textbf{120.1} \\
Phi-4 14B & No & 8.81$\times$ & 118.9 \\
Llama 4 Scout & No & 8.79$\times$ & 118.7 \\
OLMo 2 32B & Yes & 8.76$\times$ & 118.3 \\
DBRX & No & 8.71$\times$ & 117.7 \\
LLaMA 3 8B & No & 8.70$\times$ & 117.4 \\
Qwen 3 8B & Yes & 8.70$\times$ & 117.2 \\
OLMo 2 7B & Yes & 8.68$\times$ & 117.2 \\
Mixtral 8x22B & No & 8.66$\times$ & 117.0 \\
InternLM 3 8B & No & 8.66$\times$ & 116.8 \\
LLaMA 3 70B & No & 8.60$\times$ & 116.2 \\
Mistral 7B & No & 8.58$\times$ & 115.7 \\
Phi-4 mini & No & 8.50$\times$ & 114.4 \\
Qwen 3 235B-A22B & Yes & 8.32$\times$ & 112.5 \\
Gemma 4 31B global & Yes & 7.90$\times$ & 103.0 \\
Phi-3 mini & No & 7.89$\times$ & 105.1 \\
Falcon 3 10B & No & 7.58$\times$ & 101.7 \\
Falcon 3 7B & No & 7.55$\times$ & 101.2 \\
Gemma 4 E2B local & Yes & 6.49$\times$ & 79.4 \\
\bottomrule
\end{tabular}
\end{table}

Models with \texttt{head\_dim=128} cluster at 8.3--8.9$\times$; Gemma~4 models
with \texttt{head\_dim=256} land at 6.5--7.9$\times$. The explanation is launch
overhead amortization: smaller tiles have proportionally more launch overhead
relative to compute, so fusing from 4 to 2 launches saves a larger fraction.
The fusion is \textbf{not Gemma-specific}: both QK-norm and non-QK-norm
architectures benefit.


\subsection{End-to-End Model Benchmark: 26-Layer Forward Pass}
\label{sec:e2e}

\begin{table}[t]
\centering
\caption{26-layer Gemma~4 E2B forward pass (batch=1, seq\_len=512, bf16).}
\label{tab:e2e}
\small
\begin{tabular}{@{}lcccc@{}}
\toprule
Configuration & T4 (ms) & T4 speedup & A100 (ms) & A100 speedup \\
\midrule
Baseline (PyTorch eager) & 322.7 & 1.00$\times$ & 41.4 & 1.00$\times$ \\
Noeris (fused RMSNorm + QK+RoPE + GeGLU) & 274.4 & \textbf{1.18$\times$} & 29.1 & \textbf{1.43$\times$} \\
\bottomrule
\end{tabular}
\end{table}

Noeris fuses three operation classes (RMSNorm, QK-norm+RoPE, GeGLU) accounting
for ${<}20\%$ of per-layer compute, yet delivers 18\% E2E speedup on T4 and
43\% on A100. The A100 benefit is 2.4$\times$ larger than T4, confirming fusion
scales up on higher-bandwidth hardware. This is 6--21$\times$ more E2E impact
than vLLM's reported 2--3\%.


\subsection{Sliding-Window Attention: Beating cuDNN FlashAttention}
\label{sec:swa-eval}

\begin{table}[t]
\centering
\caption{Sliding-window attention vs.\ SDPA (cuDNN FlashAttention~\citep{flashattention}). Config: \texttt{BLOCK\_M=32, BLOCK\_N=32, warps=2, stages=2}.}
\label{tab:swa}
\small
\begin{tabular}{@{}cccccccc@{}}
\toprule
GPU & $N$ & $W$ & $H$ & $D$ & Noeris (ms) & SDPA (ms) & Speedup \\
\midrule
\multirow{4}{*}{T4} & 8192 & 128 & 8 & 64 & 5.68 & 7.08 & 1.25$\times$ \\
 & 8192 & 128 & 16 & 64 & 9.99 & 10.97 & 1.10$\times$ \\
 & 8192 & 64 & 8 & 64 & 3.12 & 7.08 & 2.27$\times$ \\
 & 8192 & 64 & 8 & 64 & 1.99 & 7.08 & \textbf{3.56$\times$} \\
\midrule
\multirow{4}{*}{A100} & 8192 & 128 & 8 & 64 & 1.14 & 3.42 & 3.00$\times$ \\
 & 8192 & 128 & 16 & 64 & 2.01 & 5.31 & 2.64$\times$ \\
 & 8192 & 64 & 8 & 64 & 0.55 & 3.42 & \textbf{6.24$\times$} \\
 & 8192 & 64 & 16 & 64 & 1.02 & 5.31 & 5.21$\times$ \\
\bottomrule
\end{tabular}
\end{table}

On A100, all 8 tested shapes show wins (8/8), with speedups up to
\textbf{6.24$\times$}. On T4, 4/8 shapes win. The best speedups occur at the
narrowest windows ($W$=64) where tile-pruning skips 98\% of tiles. The advantage
is specific to the narrow-window regime: $W/N$ must be small enough that
tile-pruning eliminates a large fraction of work, and \texttt{head\_dim} must be
small enough (64) that per-tile compute is cheap relative to iteration overhead.

Sliding-window attention is architecturally important: Gemma~3/4 runs 5 of 6
attention layers as local windows, Mistral uses sliding-window throughout.
Active community bounties exist in vLLM (PR~\#24390) and Axolotl (issue~\#1038).


\subsection{Cross-Hardware Transfer: Zero-Shot Config Prediction}
\label{sec:cross-hw-transfer}

\begin{table}[t]
\centering
\caption{A100 prologue fusion: 19 models, 13 families. Spearman $\rho = 0.907$ (A100 vs.\ H100).}
\label{tab:a100-19model}
\small
\begin{tabular}{@{}lcc@{}}
\toprule
Model & Fusion speedup & GB/s \\
\midrule
gemma4\_31b\_global & \textbf{9.76$\times$} & \textbf{655.0} \\
gemma4\_26b\_a4b\_global & 9.41$\times$ & 631.2 \\
gemma4\_31b\_local & 8.12$\times$ & 544.7 \\
gemma3\_local\_1024 & 7.89$\times$ & 529.2 \\
Qwen 3 32B & 5.33$\times$ & 357.5 \\
Phi-4 14B & 5.28$\times$ & 354.2 \\
Llama 4 Scout & 5.27$\times$ & 353.4 \\
DBRX & 5.22$\times$ & 350.2 \\
LLaMA 3 8B & 5.21$\times$ & 349.4 \\
OLMo 2 7B & 5.20$\times$ & 348.8 \\
Qwen 3 8B & 5.20$\times$ & 348.8 \\
Mixtral 8x22B & 5.19$\times$ & 348.1 \\
LLaMA 3 70B & 5.16$\times$ & 346.0 \\
Mistral 7B & 5.14$\times$ & 344.9 \\
Phi-4 mini & 5.09$\times$ & 341.5 \\
Qwen 3 235B-A22B & 4.98$\times$ & 334.1 \\
Phi-3 mini & 4.73$\times$ & 317.2 \\
Falcon 3 10B & 4.54$\times$ & 304.5 \\
Falcon 3 7B & 4.52$\times$ & 303.2 \\
\bottomrule
\end{tabular}
\end{table}

All 19 models achieve 3.9--9.8$\times$ fusion speedup on A100 (mean 340~GB/s).
Cross-hardware fused kernel latency rankings between A100 and H100 achieve
\textbf{Spearman $\rho = 0.907$ ($p < 10^{-7}$)}, validating zero-shot
cross-hardware configuration prediction. Prior art (TenSet,
\citealt{tenset}) requires target-device fine-tuning; Noeris achieves
strong rank correlation with \textbf{zero H100 training data}.


\subsection{Bandit Convergence: 98\% of Optimal in One Iteration}
\label{sec:bandit-convergence}

\begin{table}[t]
\centering
\caption{Bandit convergence: iteration~1 (6 configs, 12\% of search space) vs.\ exhaustive (50 configs).}
\label{tab:bandit-convergence}
\small
\begin{tabular}{@{}lccc@{}}
\toprule
Operator & Exhaustive (50 configs) & Bandit iter~1 (6 configs) & \% of optimal \\
\midrule
rmsnorm & 228.0 GB/s & 227.8 GB/s & 99.9\% \\
qk\_norm\_rope & 114.0 GB/s & 114.0 GB/s & 99.9\% \\
geglu & 248.9 GB/s & 244.5 GB/s & 98.2\% \\
\bottomrule
\end{tabular}
\end{table}

The bandit evaluates only 6 of 50 candidate configurations (12\% of the search
space) and finds $\geq$98\% of optimal throughput on the first iteration across
all three operators. Full optimality (100\%) is reached by iteration 2--6.
This represents an \textbf{8.3$\times$ reduction in GPU cost} with negligible
performance loss, confirming fully autonomous self-improvement with zero human
intervention.


\subsection{Architecture-Dependent Triton Codegen Quality}
\label{sec:codegen-quality}

The same fused kernel that delivers 12$\times$ speedup over separated Triton on
A100 produces only 0.59$\times$ vs.\ PyTorch native at the layer level on T4.
18 of 20 non-matmul operators pass correctness on T4; the issue is purely Triton
PTX quality on Turing (SM~7.5). Triton generates high-quality code for Ampere (SM~8.0) and
Hopper (SM~9.0); Turing receives less optimization attention.

\begin{table}[t]
\centering
\caption{Architecture-dependent codegen quality.}
\label{tab:codegen}
\small
\begin{tabular}{@{}lccc@{}}
\toprule
Metric & T4 (Turing) & A100 (Ampere) & H100 (Hopper) \\
\midrule
Fused vs.\ separated Triton & 6.64$\times$ & 10.2--12.9$\times$ & 10.4--11.9$\times$ \\
Fused vs.\ PyTorch native (layer) & \textbf{0.59$\times$} & $\gg$1$\times$ & $\gg$1$\times$ \\
Correctness & 18/20 & 60/60 & 60/60 \\
\bottomrule
\end{tabular}
\end{table}


% ---------------------------------------------------------------------------
% 5. RELATED WORK
% ---------------------------------------------------------------------------
\section{Related Work}
\label{sec:related}

GPU kernel optimization spans hand-tuned libraries, compiler-directed
autotuning, algebraic search, and LLM-driven code generation agents. We organize
the landscape into four groups.

\subsection{LLM-Driven Kernel Generation and Agent Loops}
\label{sec:rw-llm}

\textbf{AutoKernel}~\citep{autokernel} is the closest direct competitor:
autonomous LLM-driven edit-benchmark-keep loop on 9 Triton + 9 CUDA C++ starter
kernels, with a 5-stage correctness harness and Amdahl's-law multi-kernel
orchestrator. On H100: RMSNorm 2,788~GB/s (5.29$\times$), softmax 2,800~GB/s,
cross-entropy 2,070~GB/s; beats \texttt{torch.compile} on 12/16 configs.
AutoKernel optimizes \emph{individual kernel implementations} through 300+
experiments per kernel. Noeris operates at a different level: it discovers
cross-operator fusion patterns (QK-RMSNorm+RoPE, PLE, GeGLU) and learns
hardware-specific configurations via Thompson-sampling bandit with 50$\times$
fewer experiments (6~configs to 98\% optimal vs.\ 300+). Additional
differentiators: cross-hardware zero-shot transfer ($\rho$=0.907 fused,
$\rho$=0.967 per-op) vs.\ single-GPU; transformers + SSMs vs.\ transformers
only; fusion \emph{discovery} vs.\ per-kernel optimization. The two approaches
are complementary---AutoKernel's per-kernel LLM loop could serve as a mutation
operator within Noeris's search framework. Our curated starters already exceed
AutoKernel's best H100 figures on three of four operators.

\textbf{KernelSkill}~\citep{kernelskill} introduces multi-agent framework with
dual-level memory (long-term skills, short-term backtracking). 5.44$\times$
average on L1. The skill library is the closest analog to our database, but
stores optimization strategies rather than numerical outcomes and does not key by
\texttt{(operator, shape\_bucket, hardware)}.

\textbf{CUDA Agent}~\citep{cudaagent}: large-scale agentic RL. 100\% beat
\texttt{torch.compile} on L1--L2. Cross-run learning baked into model weights
rather than an explicit store.

\textbf{CUDA-L1}~\citep{cudal1}: contrastive RL over free-form CUDA source.
Average 3.12$\times$, peak 120$\times$. Trades tractability for expressiveness
relative to Noeris's structured parameter space.

\textbf{KernelFoundry}~\citep{kernelfoundry}: MAP-Elites quality-diversity +
LLM + meta-prompt evolution. Uses template-based parameter optimization---the
one point of contact with Noeris's templates---but templates are one tool in a
broader evolutionary repertoire.

\textbf{KernelBand}~\citep{kernelband}: hierarchical MAB, the most structurally
similar to our bandit selector. Key differences: no learned cost model, no
cross-run database, no cross-hardware transfer.


\subsection{Learned Cost Models and Traditional Autotuners}
\label{sec:rw-cost-models}

\textbf{TVM/Ansor}~\citep{ansor}: evolutionary search + XGBoost cost model over
AST features. Noeris's GBR cost model is directly inspired by Ansor but operates
on the discrete parameter grid of a fixed Triton template.

\textbf{Triton \texttt{@autotune}} is the immediate baseline. Cache discarded
across processes. Noeris generalizes it with persistent cross-run storage, LLM
proposals, and cost model ranking.

\textbf{Helion/LFBO} (Meta, 2025): Bayesian optimization for Triton autotuning.
Single hardware target; no cross-run database or transfer.


\subsection{QK-Norm + RoPE Fusion Prior Art}
\label{sec:rw-fusion}

\textbf{vLLM} \texttt{enable\_qk\_norm\_rope\_fusion}: disabled by default due to
H100 regression; 2--3\% E2E. \textbf{SGLang} \texttt{fused\_qknorm\_rope}:
limited scope. \textbf{Modular/MAX}: fused RoPE+RMSNorm targeting DeepSeek MLA,
closed-source. \textbf{Flash Normalization} (Liao et al., 2024): algebraic
optimizations. \textbf{Liger Kernel}: fuses RMSNorm and RoPE \emph{separately},
not together. \textbf{Unsloth}: 2.3$\times$ fused QK-RoPE without RMSNorm.
\textbf{FlashAttention-4}: subsumes our fusion on Blackwell B200 but does not
exist for A100/H100.

\subsection{Attention Masking and Sparse Attention}
\label{sec:rw-attention}

\textbf{FlexAttention}~\citep{flexattention}: composable block-sparse API that
can express tile-pruning. Our contribution is empirical: first published
measurement of ${>}3\times$ wins on narrow-window shapes across both T4 and A100.

\subsection{How Noeris is Actually Different}
\label{sec:rw-diff}

No individual component is new in isolation. Noeris contributes a specific
combination: \textbf{(a)}~fixed parameterized template interface bounding the
search space; \textbf{(b)}~cross-session database keyed by (operator,
shape\_bucket, hardware); \textbf{(c)}~lightweight GBR cost model trained on that
database; \textbf{(d)}~adaptive router learning per-iteration which selector to
deploy.

\begin{table}[t]
\centering
\caption{System comparison.}
\label{tab:comparison}
\small
\begin{tabular}{@{}lccccc@{}}
\toprule
System & Search & Cost model & Cross-run DB & Cross-HW & Adaptive \\
\midrule
KernelBand & Hier.\ MAB & None & No & No & No \\
KernelEvolve & LLM evol. & None & No & No & No \\
Helion/LFBO & Bayes.\ opt. & GP & No & No & No \\
TVM/Ansor & Evol.+XGB & XGB (AST) & Per-op logs & No & No \\
AutoKernel & LLM agent & None & No & No & No \\
\textbf{Noeris} & \textbf{Thompson+LLM} & \textbf{GBR} & \textbf{Yes} & \textbf{$\rho$=0.97/0.91} & \textbf{Yes} \\
\bottomrule
\end{tabular}
\end{table}


% ---------------------------------------------------------------------------
% 6. LIMITATIONS
% ---------------------------------------------------------------------------
\section{Limitations}
\label{sec:limitations}

\begin{enumerate}[leftmargin=*,itemsep=2pt]
\item \textbf{Cross-run learning ablation is negative} (\S\ref{sec:cross-run-ablation}).
The claim rests on architecture, not yet on measured gains.

\item \textbf{Cost model ablation uses a weak baseline.} The
\texttt{--no-curated} baseline is deliberately weakened.

\item \textbf{Adaptive router validated on one operator} (matmul A100 only).

\item \textbf{Attention kernel is simplified.} Does not match SDPA/FlashAttention-2/3.

\item \textbf{KernelBench subset.} 53 of 250 problems evaluated.

\item \textbf{Single hardware vendor.} Three NVIDIA GPUs; no AMD MI300.

\item \textbf{MoE router negative result.} 0.12$\times$ fusion on T4, 0.68$\times$ on A100.

\item \textbf{Metric noise.} Modal per-call variance is ${\sim}$1--3\%.

\item \textbf{Triton codegen is architecture-dependent.} 0.59$\times$ vs.\ PyTorch native
on T4 (\S\ref{sec:codegen-quality}).

\item \textbf{Per-operator surrogates ineffective.} GP/ensemble models produce
negative R$^2$ within a single operator on a single GPU.
\end{enumerate}


% ---------------------------------------------------------------------------
% 7. CONCLUSION
% ---------------------------------------------------------------------------
\section{Conclusion}
\label{sec:conclusion}

Noeris demonstrates a complete autonomous GPU kernel search pipeline with
21 parameterized Triton operators. The headline result is a fused
QK-RMSNorm+RoPE kernel achieving $10.2$--$12.9\times$ on A100 and
$10.4$--$11.9\times$ on H100, generalizing to 19 model shapes across 13
families ($6.5$--$8.9\times$ on T4, $3.9$--$9.8\times$ on A100). A 26-layer
Gemma~4 forward pass achieves $1.18\times$ E2E speedup on T4 and $1.43\times$ on
A100---6--21$\times$ more impact than vLLM's 2--3\%. A sliding-window attention
kernel beats cuDNN FlashAttention by up to $6.24\times$ on A100. Cross-hardware
fused kernel rankings transfer with Spearman $\rho = 0.907$
($p < 10^{-7}$), validating zero-shot config prediction.

The central contribution has evolved from ``a cost model improves search'' to a
complete selector-routing story with statistical validation. The cost model
improves search by +37\% on cross\_entropy; the bandit outperforms on matmul
(+134\% vs.\ +45\%); an adaptive router matches the best fixed selector within
0.5\% across 3 trials. A100-trained cost model rankings transfer to H100 with
Spearman $\rho = 0.967$.

A bandit convergence study demonstrates $\geq$98\% of exhaustive-search optimal
in a single iteration with 12\% of the search space---the strongest evidence of
autonomous self-improvement. On memory-bound kernels, curated-starter-only
results already exceed AutoKernel's published H100 numbers on RMSNorm (+120\%),
cross-entropy (+228\%), and softmax (+85\%).

All code, raw benchmark data, and reproduction scripts are available at
\url{https://github.com/PwnKit-Labs/noeris} under the MIT License.


% ---------------------------------------------------------------------------
% REFERENCES
% ---------------------------------------------------------------------------
\bibliographystyle{plainnat}

\begin{thebibliography}{20}

\bibitem[Dao et~al.(2022)]{flashattention}
Tri Dao, Daniel Y.~Fu, Stefano Ermon, Atri Rudra, and Christopher R\'{e}.
\newblock {FlashAttention}: Fast and memory-efficient exact attention with
  {IO}-awareness.
\newblock In \emph{NeurIPS}, 2022.

\bibitem[Kwon et~al.(2023)]{vllm}
Woosuk Kwon, Zhuohan Li, Siyuan Zhuang, Ying Sheng, Lianmin Zheng, Cody~Hao
  Yu, Joseph~E. Gonzalez, Hao Zhang, and Ion Stoica.
\newblock Efficient memory management for large language model serving with
  {PagedAttention}.
\newblock In \emph{SOSP}, 2023.

\bibitem[Ouyang et~al.(2025)]{kernelbench}
Anne Ouyang, Simon Guo, and Azalia Mirhoseini.
\newblock {KernelBench}: Can {LLMs} write efficient {GPU} kernels?
\newblock arXiv:2502.10517, 2025.

\bibitem[{Peking University}(2025)]{kernelband}
{Peking University}.
\newblock {KernelBand}: Hierarchical multi-armed bandit for {GPU} kernel
  optimization.
\newblock arXiv:2511.18868, 2025.

\bibitem[Chen et~al.(2021)]{tenset}
Lianmin Chen, et~al.
\newblock {TenSet}: A large-scale program performance dataset for learned
  tensor compilers.
\newblock In \emph{NeurIPS Datasets and Benchmarks Track}, 2021.

\bibitem[He et~al.(2024)]{flexattention}
Horace He, Driss Guessous, Yanbo Liang, and Joy Dong.
\newblock {FlexAttention}: The flexibility of {PyTorch} with the performance of
  {FlashAttention}.
\newblock PyTorch Blog, 2024.

\bibitem[Jaber and Jaber(2026)]{autokernel}
Eiman Jaber and Christoph Jaber.
\newblock {AutoKernel}: Autonomous {GPU} kernel optimization via iterative
  agent-driven search.
\newblock arXiv:2603.21331, 2026.

\bibitem[Sun et~al.(2026)]{kernelskill}
Sun, et~al.
\newblock {KernelSkill}: A multi-agent framework for {GPU} kernel optimization.
\newblock arXiv:2603.10085, 2026.

\bibitem[Li et~al.(2026)]{cudal1}
Li, et~al.
\newblock {CUDA-L1}: Improving {CUDA} optimization via contrastive
  reinforcement learning.
\newblock arXiv:2507.14111, ICLR 2026.

\bibitem[Dai et~al.(2026)]{cudaagent}
Dai, et~al.
\newblock {CUDA Agent}: Large-scale agentic reinforcement learning for {GPU}
  kernel optimization.
\newblock arXiv:2602.24286, 2026.

\bibitem[Wiedemann et~al.(2026)]{kernelfoundry}
Wiedemann, et~al.
\newblock {KernelFoundry}: Quality-diversity search for {GPU} kernel
  optimization.
\newblock arXiv:2603.12440, 2026.

\bibitem[Zheng et~al.(2020)]{ansor}
Lianmin Zheng, Chengfan Jia, Minmin Sun, Zhao Wu, Cody~Hao Yu, Ameer
  Haj-Ali, Yida Wang, Jun Yang, Danyang Zhuo, Koushik Sen, Joseph~E. Gonzalez,
  and Ion Stoica.
\newblock {Ansor}: Generating high-performance tensor programs for deep
  learning.
\newblock In \emph{OSDI}, 2020.

\bibitem[Tillet et~al.(2019)]{triton}
Philippe Tillet, H.~T. Kung, and David Cox.
\newblock {Triton}: An intermediate language and compiler for tiled neural
  network computations.
\newblock In \emph{MAPL@PLDI}, 2019.

\bibitem[Gu and Dao(2024)]{mamba2}
Albert Gu and Tri Dao.
\newblock Transformers are {SSMs}: Generalized models and efficient algorithms
  through structured state space duality.
\newblock arXiv:2405.21060, 2024.

\end{thebibliography}

\end{document}
